\section{Convergence Analysis}
\label{sec:convergence}

We analyze one edge or cloud learner relative to a base objective induced by complete, fresh statistics. Prototype drift compares the actual and fresh-reference feature distributions within an episode. Reference variation describes how the fresh-reference objective changes across updates.

Consider SGD-type updates $\theta_{u+1}=\theta_u-\eta g_u$ for $u=0,\ldots,T-1$, where $u$ counts valid optimizer updates. Packet ages still use logical rounds. Let $e(u)$ denote the episode that supplies update $u$, and write $Q_u=Q_{e(u)}$. We use the same convention for the reference distribution, active classes, weights, and drift. Define the population batch objective by
\begin{equation}
 F_{u,Q}(\theta)=
 \mathbb E_{\mathcal B\sim\mathsf S_u(Q)}
 [\ell_{\mathrm{base}}(\theta;\mathcal B)],
 \label{eq:tier_objective}
\end{equation}
where $\mathsf S_u$ constructs minibatches using the episode's active classes, class weights, and augmentation rule. The base loss $\ell_{\mathrm{base}}$ can be a supervised contrastive loss or a task loss. Write $F_u=F_{u,Q_u}$ for the objective under the actual feature distribution and $F_{*,u}=F_{u,Q_{*,u}}$ for its fresh-reference counterpart. Both use the same active classes and sampling rule. The drift $\drift_u$ in~\eqref{eq:prototype_drift} compares the two feature distributions using the covariances of their samplers.

\subsection{Assumptions}

Let $\mathcal F_u$ denote the information available before the stochastic draw for update $u$. It includes the current parameters, episode and synthesis state, generated pool, reference statistics, and preceding history.

\begin{assumption}[Reference regularity]
\label{asm:smooth}
Each $F_{*,u}$ is $L$-smooth and bounded below by a common $F_{\inf}$.
\end{assumption}

\begin{assumption}[Conditional update model]
\label{asm:variance}
The update direction admits the decomposition
\begin{equation}
 g_u=\nabla F_u(\theta_u)
 +b_u^{\mathrm{sys}}+b_u^{\mathrm{reg}}+\xi_u,
 \label{eq:gradient_model}
\end{equation}
where the biases are $\mathcal F_u$-measurable,
$\mathbb E[\xi_u\mid\mathcal F_u]=0$, and
$\mathbb E[\|\xi_u\|_2^2\mid\mathcal F_u]\le\sigma_g^2$.
The term $b_u^{\mathrm{reg}}$ is the conditional mean contribution from PRC, including its coefficient. The term $b_u^{\mathrm{sys}}$ captures additional conditional mean error relative to the population objective under $Q_u$, such as bias from reusing a finite pool or sampling from shuffled batches. The term $\xi_u$ accounts for the remaining stochastic variation.
\end{assumption}
Periodic model synchronization between rounds is covered whenever its parameter perturbations satisfy this update decomposition and the error bounds below.

\begin{assumption}[Distribution-sensitive gradients]
\label{asm:drift}
With active classes, weights, and batch construction held fixed,
\begin{equation}
\begin{split}
 &\|\nabla F_{u,Q}(\theta)-\nabla F_{u,Q'}(\theta)\|_2\\
 &\quad\le L_z\Bigl(
 \sum_{j\in\mathcal J_u}\pi_u^{(j)}
 W_2^2(Q^{(j)},Q'^{(j)})\Bigr)^{1/2}.
\end{split}
\label{eq:gradient_distribution_sensitivity}
\end{equation}
\end{assumption}

If the batch sampler admits a class-wise coupling with shared labels and augmentation, a batch parameter gradient that is uniformly $K_b$-Lipschitz in the joint vector of $b$ features gives $L_z=K_b\sqrt b$ for expected class proportions $\pi_u^{(j)}$. For the implemented batch-averaged supervised contrastive loss, $K_b=O(b^{-1/2})$ under the sufficient conditions in Appendix~\ref{supp:supcon_sensitivity}: a fixed positive temperature, independent per-view embeddings, normalized embeddings with bounded derivatives, and uniformly bounded first and second network derivatives. The all-pairs coupling does not change this scaling because each softmax weight is $O(1/b)$ under bounded logits. These conditions keep $L_z=O(1)$, but they do not cover arbitrary batch-dependent layers (such as active BatchNorm layers that couple views across a minibatch); Assumption~\ref{asm:drift} therefore remains explicit rather than being inferred from the loss alone.

\begin{assumption}[Additional mean error]
\label{asm:het_agg}
For nonnegative constants $H_{\mathrm{het}},H_{\mathrm{agg}}$,
$\mathbb E\|b_u^{\mathrm{sys}}\|_2^2
\le2(H_{\mathrm{het}}+H_{\mathrm{agg}})$.
This bound excludes effects already represented by $Q_u$.
\end{assumption}

\begin{assumption}[Reference variation]
\label{asm:slow}
The reference objectives satisfy the variation budget
\begin{equation}
 \mathbb E\sum_{u=1}^{T-1}\sup_\theta
 |F_{*,u}(\theta)-F_{*,u-1}(\theta)|\le V_T .
 \label{eq:reference_variation}
\end{equation}
This budget includes changes in lower-tier representations, synthesis state, active classes, weights, and sampling rules.
\end{assumption}

\subsection{Stationarity Bound}

The bound evaluates the gradients of the fresh-reference base objectives at the parameters visited by the learner. It separates prototype drift, additional mean update errors, and PRC contributions.

\begin{theorem}[Tier-wise stationarity]
\label{thm:convergence}
Under Assumptions~\ref{asm:smooth}--\ref{asm:slow}, let
$T\ge1$ and $\eta=\eta_0/\sqrt T\le1/L$, with $\eta_0>0$.
Define the finite quantities
$\Delta_F=\mathbb E F_{*,0}(\theta_0)-F_{\inf}$,
$\bar D_T=T^{-1}\sum_u\mathbb E[\drift_u]$, and
$\bar R_T=T^{-1}\sum_u\mathbb E\|b_u^{\mathrm{reg}}\|_2^2$,
where all sums run from $0$ to $T-1$. Then
\begin{equation}
\begin{split}
 \frac1T\sum_{u=0}^{T-1}
 \mathbb E\|\nabla F_{*,u}(\theta_u)\|_2^2
 \le{}&\frac{2(\Delta_F+V_T)}{\eta_0\sqrt T}+\frac{L\eta_0\sigma_g^2}{\sqrt T}
 +3L_z^2\bar D_T\\
 &+6(H_{\mathrm{het}}+H_{\mathrm{agg}})
 +3\bar R_T .
\end{split}
\label{eq:main_bound}
\end{equation}
Packet availability may be dependent.
\end{theorem}

\begin{proof}
\iffalse
First, write the update as a reference gradient plus bias and noise. Set $a_u=\nabla F_{*,u}(\theta_u)$ and
$b_u=\nabla F_u(\theta_u)-a_u+
b_u^{\mathrm{sys}}+b_u^{\mathrm{reg}}$,
so that $g_u=a_u+b_u+\xi_u$. Smoothness and conditional zero-mean noise give
\begin{equation}
\begin{split}
 \mathbb E_uF_{*,u}(\theta_{u+1})
 \le{}&F_{*,u}(\theta_u)
 -\eta\langle a_u,a_u+b_u\rangle\\
 &+\frac{L\eta^2}{2}
 \bigl(\|a_u+b_u\|_2^2+\sigma_g^2\bigr),
\end{split}
\label{eq:conditional_descent}
\end{equation}
where $\mathbb E_u$ denotes conditioning on $\mathcal F_u$.
Using
$-2\langle a,a+b\rangle=-\|a\|_2^2-\|a+b\|_2^2+\|b\|_2^2$
and $L\eta\le1$ yields
\begin{equation}
\begin{split}
 \mathbb E_uF_{*,u}(\theta_{u+1})
 \le{}&F_{*,u}(\theta_u)-\frac{\eta}{2}\|a_u\|_2^2\\
 &+\frac{\eta}{2}\|b_u\|_2^2
 +\frac{L\eta^2}{2}\sigma_g^2.
\end{split}
\label{eq:biased_descent}
\end{equation}

Next, the three-term squared-norm inequality and Assumptions~\ref{asm:drift} and~\ref{asm:het_agg} imply
\begin{equation}
\begin{split}
 \mathbb E\|b_u\|_2^2
 \le{}&3L_z^2\mathbb E[\drift_u]
 +6(H_{\mathrm{het}}+H_{\mathrm{agg}})+3\mathbb E\|b_u^{\mathrm{reg}}\|_2^2 .
\end{split}
\label{eq:total_bias_bound}
\end{equation}
This step requires no independence among the bias components.

Finally, account for changes in the reference objective through the identity
\begin{equation}
\begin{split}
 &\sum_{u=0}^{T-1}
 [F_{*,u}(\theta_u)-F_{*,u}(\theta_{u+1})]=F_{*,0}(\theta_0)-F_{*,T-1}(\theta_T)\\
 &\quad+\sum_{u=1}^{T-1}
 [F_{*,u}(\theta_u)-F_{*,u-1}(\theta_u)].
\end{split}
\label{eq:changing_reference_telescope}
\end{equation}
The expected right-hand side is at most $\Delta_F+V_T$ by the common lower bound and Assumption~\ref{asm:slow}. Summing~\eqref{eq:biased_descent}, taking expectations, and dividing by $\eta T/2$ gives
\begin{equation}
\begin{split}
 \frac1T\sum_{u=0}^{T-1}\mathbb E\|a_u\|_2^2
 \le{}&\frac{2(\Delta_F+V_T)}{\eta T}
 +L\eta\sigma_g^2\\
 &+\frac1T\sum_{u=0}^{T-1}\mathbb E\|b_u\|_2^2 .
\end{split}
\label{eq:average_descent}
\end{equation}
Substituting~\eqref{eq:total_bias_bound} and $\eta=\eta_0/\sqrt T$ proves~\eqref{eq:main_bound}.
\fi
The detailed descent and changing-reference argument is given in Appendix~\ref{supp:analysis}.
\end{proof}

\begin{corollary}[Fresh-prototype limit]
\label{cor:small_perturb}
Under Theorem~\ref{thm:convergence}, if
$\bar D_T=\bar R_T=H_{\mathrm{het}}=H_{\mathrm{agg}}=0$
and $V_T=O(1)$, the average squared reference-gradient norm is $O(T^{-1/2})$.
\end{corollary}
\begin{proof}
Only the first two terms of~\eqref{eq:main_bound} remain.
\end{proof}

\begin{remark}[Controlling the drift residual]
\label{rem:drift_residual}
The drift residual $3L_z^2\bar D_T$ links the bound to the mechanisms of Section~\ref{sec:formulation}. With fresh current packets, Proposition~\ref{prop:replay_drift} makes the replay contribution vanish linearly as $\alpha\to1$. Proposition~\ref{prop:masking_coverage} shows that masking enters mainly through coverage when classes are held by few sources, and Lemma~\ref{prop:drift_decomp} below makes the age contribution quadratic in the maximum contribution age $A$. The residual does not include the gaps $\varepsilon_{\mathrm{ref}}$ and $\varepsilon_{\mathrm{rep}}$ in~\eqref{eq:surrogate_hierarchy}, because the reference objective already incorporates them.
\end{remark}

The optimization and noise terms in~\eqref{eq:main_bound} scale as $O((1+V_T)/\sqrt T)$ and reduce to $O(T^{-1/2})$ when $V_T=O(1)$. For the right-hand side to vanish, the reference variation must satisfy $V_T=o(\sqrt T)$ and the residual terms for drift, mean update error, and PRC must also vanish. Because Theorem~\ref{thm:convergence} evaluates gradients conditioned on the represented classes $\mathcal J_u$, replay, masking, and PRC appear in~\eqref{eq:main_bound} solely through their within-$\mathcal J_u$ drift and regularization residuals ($3L_z^2\bar D_T$ and $3\bar R_T$); their stabilizing benefits---preserving class coverage $\kappa_u$ against missing classes $j\notin\mathcal J_u$ and smoothing cross-round reference jumps $V_T$---operate outside this fixed-$\mathcal J_u$ residual decomposition. Finally, the constant step size $\eta=\eta_0/\sqrt T$ uses a known horizon $T$ for algebraic simplicity (an anytime schedule $\eta_u=\eta_0/\sqrt{u+1}$ yields the analogous $O((1+V_T+\log T)/\sqrt T)$ scaling; see Section~\ref{sec:discussion} for the relation to the Adam experiments).

\subsection{Conditional Age Decomposition}

To relate drift to prototype age, we compare each synthetic distribution with one built by refreshing the moments of its contributing sources. This comparison requires conditions on how source statistics change over time and how these changes affect synthesis. We represent a Gaussian by its mean and positive semidefinite covariance square root $(\mu,\Sigma^{1/2})$. The norm for such pairs is $\|(v,B)\|_\oplus^2=\|v\|_2^2+\|B\|_F^2$.

\begin{assumption}[Contribution ages and pipeline sensitivity]
\label{asm:staleness}
Every underlying statistic contributing to synthesis, including replayed and accumulated contributions, has age at most
$A=\stale_{\max}+A_{\mathcal M}$ logical rounds (accounting for maximum arrival staleness $\stale_{\max}$ plus bank retention $A_{\mathcal M}$).
For each source, consecutive-round mean/covariance-square-root pairs differ by at most $\delta$ in $\|\cdot\|_\oplus$.
Holding source coverage, support counts, replay and reconstruction coefficients, realized aggregation weights, and synthesis hyperparameters fixed, perturbing every contributing source pair by at most $\epsilon$ changes each final generated pair by at most $L_{\mathrm{agg}}\epsilon$. This sensitivity condition includes covariance estimation, numerical regularization, and fallback.
\end{assumption}

Let $\widetilde Q_u^{(j)}$ denote the intermediate feature law obtained by refreshing the contributing moments while keeping the other quantities in Assumption~\ref{asm:staleness} fixed. Its remaining discrepancy from the complete fresh reference is
\begin{equation}
 \varepsilon_{\mathrm{pipe},u}
 =\sum_{j\in\mathcal J_u}\pi_u^{(j)}
 W_2^2(\widetilde Q_u^{(j)},Q_{*,u}^{(j)}).
 \label{eq:coverage_error}
\end{equation}
The law retains the episode's source coverage and other fixed quantities. The fresh reference uses complete statistics and recomputes its input-dependent synthesis quantities. The residual compares these laws over the represented classes.

\begin{lemma}[Age-dependent drift]
\label{prop:drift_decomp}
Under Assumption~\ref{asm:staleness},
\begin{equation}
 \mathbb E[\drift_u]\le
 2L_{\mathrm{agg}}^2A^2\delta^2
 +2\mathbb E[\varepsilon_{\mathrm{pipe},u}].
 \label{eq:drift_bound}
\end{equation}
\end{lemma}
\begin{proof}
\iffalse
For Gaussian laws with pairs $(\mu,B)$ and $(\mu',B')$, coupling through the same $\epsilon\sim\mathcal N(0,I_d)$ gives
$W_2^2(G,G')\le\mathbb E\|(\mu-\mu')+(B-B')\epsilon\|_2^2
=\|(\mu-\mu',B-B')\|_\oplus^2$.
Telescoping source motion over at most $A$ rounds and applying pipeline sensitivity therefore yields
$W_2^2(Q_u^{(j)},\widetilde Q_u^{(j)})
\le L_{\mathrm{agg}}^2A^2\delta^2$.
The squared triangle inequality bounds each class drift by
$2W_2^2(Q_u^{(j)},\widetilde Q_u^{(j)})
+2W_2^2(\widetilde Q_u^{(j)},Q_{*,u}^{(j)})$.
Weighting over classes and taking expectations proves the result.
\fi
The Gaussian coupling and age-decomposition algebra are given in Appendix~\ref{supp:analysis}.
\end{proof}

An untruncated accumulator can retain contributions from expired records, so applying Lemma~\ref{prop:drift_decomp} requires a separate age bound for all contributing statistics and the sensitivity condition in Assumption~\ref{asm:staleness}. Theorem~\ref{thm:convergence} uses the realized drift directly and does not require this decomposition. The bound does not cover class coverage and update frequency.